\section{Power Management and Full Pipeline Integration}
\label{sec:cs10-11}

The final two case studies close the loop: power management between radar frames, and end-to-end hardware-in-the-loop verification of the complete pipeline.

\subsection{CS10 --- Power Management with WFI}

Between radar frames, the coprocessor enters WFI sleep to reduce power consumption.

\textbf{Lost-wakeup race:} The check-then-sleep sequence (\texttt{if (no\_work) \_\_asm\_\_("wfi")}) has a race: an interrupt can fire between the check and the sleep, losing the wakeup. The fix wraps the check-and-sleep in a PRIMASK critical section (\texttt{cpsid i / wfi / cpsie i}). WFI wakes on pending interrupt regardless of PRIMASK; the ISR runs after interrupts are re-enabled.

\textbf{Sleep mode vs stop mode:} The \texttt{SLEEPDEEP} bit in SCB\_CR was initially set, stopping all clocks including TIM2. The fix uses Sleep mode (SLEEPDEEP=0) where peripherals remain clocked.

\subsection{CS11 --- Full Pipeline Integration and HIL Verification}

The complete pipeline is verified end-to-end in Renode with real MAVLink data.

\textbf{Parser multi-packet loss:} The initial parser processed all packets in one call but only returned the last packet's data, losing earlier packets. The fix processes one packet per call and saves the read position for the next invocation.

\textbf{Ring buffer overflow on burst:} The parser pushed to the same \texttt{state} variable used by the overflow pop, overwriting the new state with the popped state. The fix uses a separate \texttt{dummy} variable for the pop.

\textbf{FPU enable:} The firmware initially hung on the first floating-point store because the Cortex-M7 FPU was not enabled. The fix adds FPU enable to Reset\_Handler.

\textbf{Hardware-in-the-loop results:} 50 frames of pymavlink-generated MAVLink v2 data (ATTITUDE + LOCAL\_POSITION\_NED at 100\,Hz, frame-sync at 10\,Hz) are injected via AXI SRAM backdoors. The pipeline produces georeferenced output (lat$\approx$13, lon$\approx$80, alt$\approx$0\,m) consistent with the reference position and target geometry.

Complete narrative: \texttt{docs/engineering-notebook/010-power-wfi.md} and \texttt{011-full-pipeline.md}.